\begin{figure}[h!]
    \centering
    \caption{Propuesta para el ciclo de aprobación del PGN}
    \label{fig:flujoPGN}
    \vspace{0.3cm}
\begin{tikzpicture}[
    node distance = 0.8cm, % menor distancia vertical
    block/.style = {
        rectangle,
        draw,
        text width=12cm, % más ancho
        minimum height=0.6cm, % más bajo
        align=center
    },
    arrow/.style = {-Stealth, thick}
]

% Nodos
\node[block] (1) {1. CARF valida proyección de ingresos y OATP valida inflexibilidades legales de gasto};
\node[block, below=of 1] (2) {2. Hacienda presenta PGN con todas las obligaciones legales de gasto};
\node[block, below=of 2] (3) {3. Si los ingresos proyectados son insuficientes, Hacienda anuncia riesgo de incumplir con la regla fiscal.};
\node[block, below=of 3] (4) {4. Se presenta Ley de Financiamiento en caso de que proyecciones de ingresos y obligaciones legales de gasto sean inconsistentes con la regla fiscal};
%\node[block, below=of 4] (5) {5. Si la Ley de Financiamiento no es aprobada, el Gobierno decreta los impuestos apenas suficientes para cumplir con las inflexibilidades legales.};
\node[block, below=of 4] (5) {5. Si la Ley de Financiamiento no es aprobada se declara la emergencia económica};
\node[block, below=of 5] (6) {6. El Gobierno decreta los impuestos apenas suficientes para cumplir con las inflexibilidades legales y con la regla fiscal.};

% Flechas
\draw[arrow] (1) -- (2);
\draw[arrow] (2) -- (3);
\draw[arrow] (3) -- (4);
\draw[arrow] (4) -- (5);
\draw[arrow] (5) -- (6);
\end{tikzpicture}
\end{figure}
